\lyseis{Β3.6}{Κατακορυφήν, συμπληρωματικές και παραπληρωματικές γωνίες}
\sol{1} α) Λάθος: αυτές είναι συμπληρωματικές· παραπληρωματικές είναι εκείνες με άθροισμα $180^\circ$. β) Σωστό. γ) Σωστό: $90^\circ-0^\circ=90^\circ$. δ) Σωστό: αν η γωνία είναι μεγαλύτερη από $90^\circ$, η παραπληρωματική της είναι $180^\circ$ μείον κάτι μεγαλύτερο από $90^\circ$, άρα μικρότερη από $90^\circ$ (οξεία). ε) Λάθος: αρκεί να έχουν το ίδιο μέτρο· δεν έχουν υποχρεωτικά αντικείμενες πλευρές. στ) Σωστό.
\sol{2} $55^\circ$, $28^\circ$, $82^\circ$, $1^\circ$.
\sol{3} $50^\circ$, $90^\circ$, $175^\circ$, $69^\circ$.
\sol{4} Η κατακορυφήν της είναι $48^\circ$. Οι διπλανές της είναι παραπληρωματικές της: $180^\circ-48^\circ=132^\circ$, και οι δύο.
\sol{5} Κατακορυφήν: $x\hat{\mathrm{Ο}}y$ και $x'\hat{\mathrm{Ο}}y'$, καθώς και $y\hat{\mathrm{Ο}}x'$ και $y'\hat{\mathrm{Ο}}x$. Παραπληρωματικές: $x\hat{\mathrm{Ο}}y$ με $y\hat{\mathrm{Ο}}x'$, $y\hat{\mathrm{Ο}}x'$ με $x'\hat{\mathrm{Ο}}y'$, $x'\hat{\mathrm{Ο}}y'$ με $y'\hat{\mathrm{Ο}}x$ και $y'\hat{\mathrm{Ο}}x$ με $x\hat{\mathrm{Ο}}y$.
\sol{6} Ελεύθερη κατασκευή. Η συμπληρωματική είναι $90^\circ-40^\circ=50^\circ$ και η παραπληρωματική $180^\circ-40^\circ=140^\circ$.
\sol{7} α) Αμβλεία. β) Ορθή. γ) Οξεία. δ) Οξεία.
\sol{8} $\mathrm{ΑΟΔ}=30^\circ+45^\circ+50^\circ=125^\circ$. Είναι αμβλεία.
\sol{9} Οι γωνίες είναι $70^\circ$, $110^\circ$, $70^\circ$, $110^\circ$. Εμφανίζονται δύο διαφορετικές τιμές.
\sol{10} $90^\circ-35^\circ=55^\circ$.
\sol{11} $90^\circ-72^\circ=18^\circ$.
\sol{12} Οι γωνίες των λαβών είναι κατακορυφήν των γωνιών των λεπίδων, άρα $40^\circ$.
\sol{13} $\frac{2}{5}\cdot90^\circ=36^\circ$. Συμπληρωματική: $90^\circ-36^\circ=54^\circ$. Παραπληρωματική: $180^\circ-36^\circ=144^\circ$.
\sol{14} Αν $\hat{\alpha}$ και $\hat{\beta}$ είναι δύο διπλανές γωνίες δύο τεμνόμενων ευθειών και $\hat{\gamma}$ η κατακορυφήν της $\hat{\alpha}$, τότε $\hat{\alpha}+\hat{\beta}=180^\circ$ και $\hat{\gamma}+\hat{\beta}=180^\circ$, άρα $\hat{\alpha}=180^\circ-\hat{\beta}=\hat{\gamma}$.
\sol{15} Όχι υποχρεωτικά. Οι δύο γωνίες της βάσης ενός ισοσκελούς τριγώνου είναι ίσες, αλλά δεν είναι κατακορυφήν, γιατί δεν έχουν κοινή κορυφή.
\sol{16} α) Η γωνία και η συμπληρωματική της έχουν άθροισμα $90^\circ$ και διαφορά $20^\circ$. Η μικρότερη (η γωνία) είναι $(90^\circ-20^\circ):2=35^\circ$. β) Η γωνία και η παραπληρωματική της έχουν άθροισμα $180^\circ$, και η παραπληρωματική είναι $40^\circ$ μικρότερη, άρα η γωνία είναι $(180^\circ+40^\circ):2=110^\circ$. γ) $x+2x=90^\circ$, $3x=90^\circ$, $x=30^\circ$. δ) $x+3x=180^\circ$, $4x=180^\circ$, $x=45^\circ$.
\sol{17} Ναι, η $45^\circ$ (γιατί $45^\circ+45^\circ=90^\circ$) και η $90^\circ$ (γιατί $90^\circ+90^\circ=180^\circ$).
\sol{18} Ελεύθερη κατασκευή. Οι κατακορυφήν γωνίες βγαίνουν ίσες (μέχρι τα σφάλματα μέτρησης) και οι διπλανές έχουν άθροισμα $180^\circ$.
\sol{19} Ελεύθερη απάντηση. Ο μαθητής πρέπει να σημειώσει το αντικείμενο και τις γωνίες που μέτρησε.
